\documentclass[11pt]{amsart}
\usepackage{geometry}                % See geometry.pdf to learn the layout options. There are lots.
\geometry{letterpaper}                   % ... or a4paper or a5paper or ... 
%\geometry{landscape}                % Activate for for rotated page geometry
%\usepackage[parfill]{parskip}    % Activate to begin paragraphs with an empty line rather than an indent
\usepackage{graphicx}
\usepackage{amssymb}
\usepackage{epstopdf}
\usepackage{color}
\usepackage{colortbl}
\DeclareGraphicsRule{.tif}{png}{.png}{`convert #1 `dirname #1`/`basename #1 .tif`.png}

\title{PS 5.1}
%\date{}                                           % Activate to display a given date or no date

\begin{document}
\maketitle
%\section{}
%\subsection{}

\noindent Una compañía distribuye N productos a distintos comercios de la ciudad. Para ello almacena en un arreglo toda la información relacionada a su mercancía:\\

\noindent • Clave. \\

\noindent • Descripción. \\

\noindent • Existencia.\\

\noindent • Mínimo a mantener en existencia. \\

\noindent •Precio unitario.\\

\noindent Escriba un diagrama de flujo que pueda llevar a cabo las siguientes operaciones:\\

\noindent a) Venta de un producto: Se deben actualizar los campos que correspondan y verificar que la nueva existencia no esté por debajo del mínimo. (Datos: clave, cantidad vendida).\\

\noindent b) Reabastecimiento de un producto: Se deben actualizar los campos que correspondan (Datos: clave, cantidad comprada).\\

\noindent c) Actualizar el precio de un producto (Datos: clave, porcentaje de aumento).\\

\noindent d) Informar sobre un producto: Se deben proporcionar todos los datos relacionados a un producto (Dato: clave).\\

\noindent Dato: PRODUCTOS[1..N] 1 < N < 150

\noindent Donde:\\

\noindent PRODUCTOS: es un arreglo unidimensional de registros. Cada registro tiene los siguientes campos:\\

\noindent CLA: Variable de tipo entero. Representa la clave del producto.\\

\noindent EXI: Variable de tipo entero. Representa la existencia del producto.\\

\noindent MIN: Variable de tipo entero. Representa el mínimo a mantener en existencia del producto con clave CLA.\\

\noindent PU: Variable de tipo real. Expresa el precio unitario del producto.\\
\end{document}   